% Options for packages loaded elsewhere
% Options for packages loaded elsewhere
\PassOptionsToPackage{unicode}{hyperref}
\PassOptionsToPackage{hyphens}{url}
%
\documentclass[
  spanish,
  letterpaper,
  a4paper,
  11pt,
  oneside,
  numbers=noenddot]{scrbook}
\usepackage{xcolor}
\usepackage{amsmath,amssymb}
\setcounter{secnumdepth}{5}
\usepackage{iftex}
\ifPDFTeX
  \usepackage[T1]{fontenc}
  \usepackage[utf8]{inputenc}
  \usepackage{textcomp} % provide euro and other symbols
\else % if luatex or xetex
  \usepackage{unicode-math} % this also loads fontspec
  \defaultfontfeatures{Scale=MatchLowercase}
  \defaultfontfeatures[\rmfamily]{Ligatures=TeX,Scale=1}
\fi
\usepackage{lmodern}
\ifPDFTeX\else
  % xetex/luatex font selection
  \setmainfont[]{Libertinus Serif}
  \setmonofont[Scale=0.85]{DejaVu Sans Mono}
  \setmathfont[]{Libertinus Math}
\fi
% Use upquote if available, for straight quotes in verbatim environments
\IfFileExists{upquote.sty}{\usepackage{upquote}}{}
\IfFileExists{microtype.sty}{% use microtype if available
  \usepackage[]{microtype}
  \UseMicrotypeSet[protrusion]{basicmath} % disable protrusion for tt fonts
}{}
\makeatletter
\@ifundefined{KOMAClassName}{% if non-KOMA class
  \IfFileExists{parskip.sty}{%
    \usepackage{parskip}
  }{% else
    \setlength{\parindent}{0pt}
    \setlength{\parskip}{6pt plus 2pt minus 1pt}}
}{% if KOMA class
  \KOMAoptions{parskip=half}}
\makeatother
% Make \paragraph and \subparagraph free-standing
\makeatletter
\ifx\paragraph\undefined\else
  \let\oldparagraph\paragraph
  \renewcommand{\paragraph}{
    \@ifstar
      \xxxParagraphStar
      \xxxParagraphNoStar
  }
  \newcommand{\xxxParagraphStar}[1]{\oldparagraph*{#1}\mbox{}}
  \newcommand{\xxxParagraphNoStar}[1]{\oldparagraph{#1}\mbox{}}
\fi
\ifx\subparagraph\undefined\else
  \let\oldsubparagraph\subparagraph
  \renewcommand{\subparagraph}{
    \@ifstar
      \xxxSubParagraphStar
      \xxxSubParagraphNoStar
  }
  \newcommand{\xxxSubParagraphStar}[1]{\oldsubparagraph*{#1}\mbox{}}
  \newcommand{\xxxSubParagraphNoStar}[1]{\oldsubparagraph{#1}\mbox{}}
\fi
\makeatother

\usepackage{color}
\usepackage{fancyvrb}
\newcommand{\VerbBar}{|}
\newcommand{\VERB}{\Verb[commandchars=\\\{\}]}
\DefineVerbatimEnvironment{Highlighting}{Verbatim}{commandchars=\\\{\}}
% Add ',fontsize=\small' for more characters per line
\usepackage{framed}
\definecolor{shadecolor}{RGB}{241,243,245}
\newenvironment{Shaded}{\begin{snugshade}}{\end{snugshade}}
\newcommand{\AlertTok}[1]{\textcolor[rgb]{0.68,0.00,0.00}{#1}}
\newcommand{\AnnotationTok}[1]{\textcolor[rgb]{0.37,0.37,0.37}{#1}}
\newcommand{\AttributeTok}[1]{\textcolor[rgb]{0.40,0.45,0.13}{#1}}
\newcommand{\BaseNTok}[1]{\textcolor[rgb]{0.68,0.00,0.00}{#1}}
\newcommand{\BuiltInTok}[1]{\textcolor[rgb]{0.00,0.23,0.31}{#1}}
\newcommand{\CharTok}[1]{\textcolor[rgb]{0.13,0.47,0.30}{#1}}
\newcommand{\CommentTok}[1]{\textcolor[rgb]{0.37,0.37,0.37}{#1}}
\newcommand{\CommentVarTok}[1]{\textcolor[rgb]{0.37,0.37,0.37}{\textit{#1}}}
\newcommand{\ConstantTok}[1]{\textcolor[rgb]{0.56,0.35,0.01}{#1}}
\newcommand{\ControlFlowTok}[1]{\textcolor[rgb]{0.00,0.23,0.31}{\textbf{#1}}}
\newcommand{\DataTypeTok}[1]{\textcolor[rgb]{0.68,0.00,0.00}{#1}}
\newcommand{\DecValTok}[1]{\textcolor[rgb]{0.68,0.00,0.00}{#1}}
\newcommand{\DocumentationTok}[1]{\textcolor[rgb]{0.37,0.37,0.37}{\textit{#1}}}
\newcommand{\ErrorTok}[1]{\textcolor[rgb]{0.68,0.00,0.00}{#1}}
\newcommand{\ExtensionTok}[1]{\textcolor[rgb]{0.00,0.23,0.31}{#1}}
\newcommand{\FloatTok}[1]{\textcolor[rgb]{0.68,0.00,0.00}{#1}}
\newcommand{\FunctionTok}[1]{\textcolor[rgb]{0.28,0.35,0.67}{#1}}
\newcommand{\ImportTok}[1]{\textcolor[rgb]{0.00,0.46,0.62}{#1}}
\newcommand{\InformationTok}[1]{\textcolor[rgb]{0.37,0.37,0.37}{#1}}
\newcommand{\KeywordTok}[1]{\textcolor[rgb]{0.00,0.23,0.31}{\textbf{#1}}}
\newcommand{\NormalTok}[1]{\textcolor[rgb]{0.00,0.23,0.31}{#1}}
\newcommand{\OperatorTok}[1]{\textcolor[rgb]{0.37,0.37,0.37}{#1}}
\newcommand{\OtherTok}[1]{\textcolor[rgb]{0.00,0.23,0.31}{#1}}
\newcommand{\PreprocessorTok}[1]{\textcolor[rgb]{0.68,0.00,0.00}{#1}}
\newcommand{\RegionMarkerTok}[1]{\textcolor[rgb]{0.00,0.23,0.31}{#1}}
\newcommand{\SpecialCharTok}[1]{\textcolor[rgb]{0.37,0.37,0.37}{#1}}
\newcommand{\SpecialStringTok}[1]{\textcolor[rgb]{0.13,0.47,0.30}{#1}}
\newcommand{\StringTok}[1]{\textcolor[rgb]{0.13,0.47,0.30}{#1}}
\newcommand{\VariableTok}[1]{\textcolor[rgb]{0.07,0.07,0.07}{#1}}
\newcommand{\VerbatimStringTok}[1]{\textcolor[rgb]{0.13,0.47,0.30}{#1}}
\newcommand{\WarningTok}[1]{\textcolor[rgb]{0.37,0.37,0.37}{\textit{#1}}}

\usepackage{longtable,booktabs,array}
\newcounter{none} % for unnumbered tables
\usepackage{calc} % for calculating minipage widths
% Correct order of tables after \paragraph or \subparagraph
\usepackage{etoolbox}
\makeatletter
\patchcmd\longtable{\par}{\if@noskipsec\mbox{}\fi\par}{}{}
\makeatother
% Allow footnotes in longtable head/foot
\IfFileExists{footnotehyper.sty}{\usepackage{footnotehyper}}{\usepackage{footnote}}
\makesavenoteenv{longtable}
\usepackage{graphicx}
\makeatletter
\newsavebox\pandoc@box
\newcommand*\pandocbounded[1]{% scales image to fit in text height/width
  \sbox\pandoc@box{#1}%
  \Gscale@div\@tempa{\textheight}{\dimexpr\ht\pandoc@box+\dp\pandoc@box\relax}%
  \Gscale@div\@tempb{\linewidth}{\wd\pandoc@box}%
  \ifdim\@tempb\p@<\@tempa\p@\let\@tempa\@tempb\fi% select the smaller of both
  \ifdim\@tempa\p@<\p@\scalebox{\@tempa}{\usebox\pandoc@box}%
  \else\usebox{\pandoc@box}%
  \fi%
}
% Set default figure placement to htbp
\def\fps@figure{htbp}
\makeatother


% definitions for citeproc citations
\NewDocumentCommand\citeproctext{}{}
\NewDocumentCommand\citeproc{mm}{%
  \begingroup\def\citeproctext{#2}\cite{#1}\endgroup}
\makeatletter
 % allow citations to break across lines
 \let\@cite@ofmt\@firstofone
 % avoid brackets around text for \cite:
 \def\@biblabel#1{}
 \def\@cite#1#2{{#1\if@tempswa , #2\fi}}
\makeatother
\newlength{\cslhangindent}
\setlength{\cslhangindent}{1.5em}
\newlength{\csllabelwidth}
\setlength{\csllabelwidth}{3em}
\newenvironment{CSLReferences}[2] % #1 hanging-indent, #2 entry-spacing
 {\begin{list}{}{%
  \setlength{\itemindent}{0pt}
  \setlength{\leftmargin}{0pt}
  \setlength{\parsep}{0pt}
  % turn on hanging indent if param 1 is 1
  \ifodd #1
   \setlength{\leftmargin}{\cslhangindent}
   \setlength{\itemindent}{-1\cslhangindent}
  \fi
  % set entry spacing
  \setlength{\itemsep}{#2\baselineskip}}}
 {\end{list}}
\usepackage{calc}
\newcommand{\CSLBlock}[1]{\hfill\break\parbox[t]{\linewidth}{\strut\ignorespaces#1\strut}}
\newcommand{\CSLLeftMargin}[1]{\parbox[t]{\csllabelwidth}{\strut#1\strut}}
\newcommand{\CSLRightInline}[1]{\parbox[t]{\linewidth - \csllabelwidth}{\strut#1\strut}}
\newcommand{\CSLIndent}[1]{\hspace{\cslhangindent}#1}

\ifLuaTeX
\usepackage[bidi=basic,shorthands=off]{babel}
\else
\usepackage[bidi=default,shorthands=off]{babel}
\fi
\ifPDFTeX
\else
\babelfont{rm}[]{Libertinus Serif}
\fi
\ifLuaTeX
  \usepackage{selnolig} % disable illegal ligatures
\fi


\setlength{\emergencystretch}{3em} % prevent overfull lines

\providecommand{\tightlist}{%
  \setlength{\itemsep}{0pt}\setlength{\parskip}{0pt}}



 


% Preámbulo de prueba para la salida LaTeX (LuaLaTeX).
\usepackage{imakeidx}
\makeindex[title=Índice analítico, columns=2, intoc]
\addtokomafont{disposition}{\rmfamily}
\KOMAoptions{numbers=noenddot}
% Entornos propios (filtros/entornos.lua): sin numerar, cuerpo en redonda, sangría a la izquierda.
% El filtro escribe el título en línea al comienzo del primer párrafo (\atlasnotatitulo).
% No se usa amsthm con nota opcional: provoca un error de babel (bidi) en LuaLaTeX con TeX Live 2023.
\newenvironment{atlasnota}{\par\medskip\begin{list}{}{\leftmargin=1.2em\rightmargin=0pt}\item\relax}{\end{list}\medskip}
\newcommand{\atlasnotatitulo}[1]{\textbf{#1.}\enspace}
% Números de sección sin punto final (el español de babel lo añade por omisión).
\renewcommand*{\autodot}{}
% Corte de líneas largas en los bloques de código.
\usepackage{fvextra}
\DefineVerbatimEnvironment{Highlighting}{Verbatim}{breaklines,breakanywhere,commandchars=\\\{\}}
\makeatletter
\@ifpackageloaded{tcolorbox}{}{\usepackage[skins,breakable]{tcolorbox}}
\@ifpackageloaded{fontawesome5}{}{\usepackage{fontawesome5}}
\definecolor{quarto-callout-color}{HTML}{909090}
\definecolor{quarto-callout-note-color}{HTML}{0758E5}
\definecolor{quarto-callout-important-color}{HTML}{CC1914}
\definecolor{quarto-callout-warning-color}{HTML}{EB9113}
\definecolor{quarto-callout-tip-color}{HTML}{00A047}
\definecolor{quarto-callout-caution-color}{HTML}{FC5300}
\definecolor{quarto-callout-color-frame}{HTML}{acacac}
\definecolor{quarto-callout-note-color-frame}{HTML}{4582ec}
\definecolor{quarto-callout-important-color-frame}{HTML}{d9534f}
\definecolor{quarto-callout-warning-color-frame}{HTML}{f0ad4e}
\definecolor{quarto-callout-tip-color-frame}{HTML}{02b875}
\definecolor{quarto-callout-caution-color-frame}{HTML}{fd7e14}
\makeatother
\makeatletter
\@ifpackageloaded{bookmark}{}{\usepackage{bookmark}}
\makeatother
\makeatletter
\@ifpackageloaded{caption}{}{\usepackage{caption}}
\AtBeginDocument{%
\ifdefined\contentsname
  \renewcommand*\contentsname{Tabla de contenidos}
\else
  \newcommand\contentsname{Tabla de contenidos}
\fi
\ifdefined\listfigurename
  \renewcommand*\listfigurename{Listado de Figuras}
\else
  \newcommand\listfigurename{Listado de Figuras}
\fi
\ifdefined\listtablename
  \renewcommand*\listtablename{Listado de Tablas}
\else
  \newcommand\listtablename{Listado de Tablas}
\fi
\ifdefined\figurename
  \renewcommand*\figurename{Figura}
\else
  \newcommand\figurename{Figura}
\fi
\ifdefined\tablename
  \renewcommand*\tablename{Tabla}
\else
  \newcommand\tablename{Tabla}
\fi
}
\@ifpackageloaded{float}{}{\usepackage{float}}
\floatstyle{ruled}
\@ifundefined{c@chapter}{\newfloat{codelisting}{h}{lop}}{\newfloat{codelisting}{h}{lop}[chapter]}
\floatname{codelisting}{Listado}
\newcommand*\listoflistings{\listof{codelisting}{Listado de Listados}}
\usepackage{amsthm}
\theoremstyle{definition}
\newtheorem{definition}{Definición}[chapter]
\theoremstyle{plain}
\newtheorem{corollary}{Corolario}[chapter]
\theoremstyle{definition}
\newtheorem{exercise}{Ejercicio}[chapter]
\theoremstyle{definition}
\newtheorem{example}{Ejemplo}[chapter]
\theoremstyle{plain}
\newtheorem{theorem}{Teorema}[chapter]
\theoremstyle{remark}
\AtBeginDocument{\renewcommand*{\proofname}{Demostración}}
\newtheorem*{remark}{Observación}
\newtheorem*{solution}{Solución}
\newtheorem{refremark}{Observación}[chapter]
\newtheorem{refsolution}{Solución}[chapter]
\makeatother
\makeatletter
\makeatother
\makeatletter
\@ifpackageloaded{caption}{}{\usepackage{caption}}
\@ifpackageloaded{subcaption}{}{\usepackage{subcaption}}
\makeatother
\usepackage{bookmark}
\IfFileExists{xurl.sty}{\usepackage{xurl}}{} % add URL line breaks if available
\urlstyle{same}
\hypersetup{
  pdftitle={Muestra de composición},
  pdfauthor={{[}nombre del director de la obra{]}},
  pdflang={es},
  hidelinks,
  pdfcreator={LaTeX via pandoc}}


\title{Muestra de composición}
\usepackage{etoolbox}
\makeatletter
\providecommand{\subtitle}[1]{% add subtitle to \maketitle
  \apptocmd{\@title}{\par {\large #1 \par}}{}{}
}
\makeatother
\subtitle{Atlas Quant · Prueba de la Fase 2}
\author{{[}nombre del director de la obra{]}}
\date{}
\begin{document}
\frontmatter
\maketitle

\renewcommand*\contentsname{Tabla de contenidos}
{
\setcounter{tocdepth}{1}
\tableofcontents
}

\mainmatter
\bookmarksetup{startatroot}

\chapter*{Créditos}\label{cruxe9ditos}
\addcontentsline{toc}{chapter}{Créditos}

\markboth{Créditos}{Créditos}

Dirección editorial, alcance y aprobación de contenidos: {[}nombre del
director de la obra{]}.\\
Redacción con asistencia de inteligencia artificial (Claude, de
Anthropic) bajo dicha dirección.

\bookmarksetup{startatroot}

\chapter{Logaritmos}\label{sec-logaritmos}

Este capítulo de muestra existe solo para probar la composición.
Reproduce, en pequeño, los elementos de un capítulo de la serie I:
definiciones, teoremas con demostración, un ejemplo financiero con
cálculo verificado por código, una tabla, una figura y ejercicios con
solución. Los textos de contraste para este tema son Spivak
(\citeproc{ref-spivak2008}{2008}) y Apostol
(\citeproc{ref-apostol1967}{1967}); la aplicación financiera sigue a
Kellison (\citeproc{ref-kellison2009}{2009}).

\section{Definición y primeras propiedades}\label{sec-log-definicion}

Sea \(b\) un número real positivo distinto de \(1\). La función
exponencial de base \(b\) es estrictamente monótona y toma todos los
valores positivos, de modo que para cada \(x > 0\) existe exactamente un
número real \(y\) tal que \(b^{y} = x\). Ese número recibe un nombre
propio.

\begin{definition}[Logaritmo]\protect\hypertarget{def-logaritmo}{}\label{def-logaritmo}

Sean \(b > 0\), \(b \neq 1\) y \(x > 0\). El
logaritmo\index{logaritmo@logaritmo} de \(x\) en base \(b\), denotado
\(\log_b x\), es el único número real \(y\) que cumple
\begin{equation}\protect\phantomsection\label{eq-def-log}{
b^{y} = x .
}\end{equation}

\end{definition}

La definición traduce una pregunta sobre exponentes en una cantidad con
nombre: \(\log_{2} 1024 = 10\) porque \(2^{10} = 1024\), y
\(\log_{10} 0{,}001 = -3\) porque \(10^{-3} = 0{,}001\). Cuando la base
es el número \(e\) se escribe \(\ln x\) y se habla de logaritmo
natural\index{logaritmo natural@logaritmo natural}.\footnote{La
  existencia de \(e\) y de la función exponencial real se establece en
  L2 c.~7; aquí se usan sin demostración porque el capítulo es una
  muestra tipográfica.}

\begin{theorem}[Logaritmo de un
producto]\protect\hypertarget{thm-log-producto}{}\label{thm-log-producto}

Para todo \(b > 0\) con \(b \neq 1\) y todos \(x, y > 0\),
\begin{equation}\protect\phantomsection\label{eq-log-producto}{
\log_b (xy) = \log_b x + \log_b y .
}\end{equation}

\end{theorem}

\begin{proof}
Sean \(u = \log_b x\) y \(v = \log_b y\). Por la definición
\ref{eq-def-log}, \(b^{u} = x\) y \(b^{v} = y\). La ley de los
exponentes da \(b^{u+v} = b^{u} b^{v} = xy\). Como \(\log_b(xy)\) es el
único número cuyo exponencial de base \(b\) es \(xy\), se concluye que
\(\log_b(xy) = u + v\).
\end{proof}

\begin{corollary}[]\protect\hypertarget{cor-log-potencia}{}\label{cor-log-potencia}

Para todo entero \(n \geq 1\) y todo \(x > 0\),
\(\log_b (x^{n}) = n \log_b x\).

\end{corollary}

\begin{proof}
Por inducción sobre \(n\). El caso \(n = 1\) es inmediato. Si la
igualdad vale para \(n\), el Teorema~\ref{thm-log-producto} aplicado a
\(x^{n}\) y \(x\) da
\(\log_b(x^{n+1}) = \log_b(x^{n}) + \log_b x = (n+1) \log_b x\).
\end{proof}

\begin{tcolorbox}[enhanced jigsaw, arc=.35mm, bottomrule=.15mm, bottomtitle=1mm, breakable, colback=white, colbacktitle=quarto-callout-warning-color!10!white, colframe=quarto-callout-warning-color-frame, coltitle=black, left=2mm, leftrule=.75mm, opacityback=0, opacitybacktitle=0.6, rightrule=.15mm, title=\textcolor{quarto-callout-warning-color}{\faExclamationTriangle}\hspace{0.5em}{Error frecuente}, titlerule=0mm, toprule=.15mm, toptitle=1mm]

El logaritmo de una suma no es la suma de los logaritmos. Por ejemplo,
\(\log_{10}(10 + 10) = \log_{10} 20 \approx 1{,}301\), mientras que
\(\log_{10} 10 + \log_{10} 10 = 2\).

\end{tcolorbox}

\begin{theorem}[Cambio de
base]\protect\hypertarget{thm-cambio-base}{}\label{thm-cambio-base}

Si \(a, b > 0\), \(a \neq 1\), \(b \neq 1\) y \(x > 0\), entonces
\begin{equation}\protect\phantomsection\label{eq-cambio-base}{
\log_b x = \frac{\log_a x}{\log_a b} .
}\end{equation}

\end{theorem}

\begin{proof}
Sea \(y = \log_b x\), de modo que \(b^{y} = x\). Aplicando \(\log_a\) y
el Corolario~\ref{cor-log-potencia} (extendido a exponentes reales en
L2), \(y \log_a b = \log_a x\). Como \(b \neq 1\), \(\log_a b \neq 0\) y
basta dividir.
\end{proof}

\section{Aplicación: tiempo de duplicación de un
capital}\label{sec-duplicacion}

Un capital de \(C_0 = 15\,000\) dólares invertido a una tasa efectiva
anual \(r\) vale \(C_0 (1 + r)^{t}\) al cabo de \(t\) años. El capital
se duplica cuando \((1 + r)^{t} = 2\), es decir, por la
Ecuación~\ref{eq-def-log} y la Ecuación~\ref{eq-cambio-base},
\begin{equation}\protect\phantomsection\label{eq-duplicacion}{
t^{*} = \log_{1+r} 2 = \frac{\ln 2}{\ln(1 + r)} .
}\end{equation}

\begin{example}[]\protect\hypertarget{exm-duplicacion}{}\label{exm-duplicacion}

Con \(r = 5\) \%, la Ecuación~\ref{eq-duplicacion} da
\(t^{*} = \ln 2 / \ln 1{,}05\). El valor se calcula en la celda
siguiente, que se ejecuta al componer el documento.

\end{example}

\protect\phantomsection\label{duplicacion}
\begin{Shaded}
\begin{Highlighting}[]
\ImportTok{import}\NormalTok{ math}

\KeywordTok{def}\NormalTok{ coma(x, decimales):}
    \CommentTok{"""Formatea un número con coma decimal y espacio fino como separador de millares."""}
\NormalTok{    texto }\OperatorTok{=} \SpecialStringTok{f"}\SpecialCharTok{\{}\NormalTok{x}\SpecialCharTok{:,.}\NormalTok{\{decimales\}f}\SpecialCharTok{\}}\SpecialStringTok{"}
    \ControlFlowTok{return}\NormalTok{ texto.replace(}\StringTok{","}\NormalTok{, }\StringTok{"}\CharTok{\textbackslash{}u202f}\StringTok{"}\NormalTok{).replace(}\StringTok{"."}\NormalTok{, }\StringTok{","}\NormalTok{).replace(}\StringTok{"{-}"}\NormalTok{, }\StringTok{"}\CharTok{\textbackslash{}u2212}\StringTok{"}\NormalTok{)}

\NormalTok{r }\OperatorTok{=} \FloatTok{0.05}
\NormalTok{t }\OperatorTok{=}\NormalTok{ math.log(}\DecValTok{2}\NormalTok{) }\OperatorTok{/}\NormalTok{ math.log(}\DecValTok{1} \OperatorTok{+}\NormalTok{ r)}
\BuiltInTok{print}\NormalTok{(}\SpecialStringTok{f"t* = }\SpecialCharTok{\{}\NormalTok{coma(t, }\DecValTok{4}\NormalTok{)}\SpecialCharTok{\}}\SpecialStringTok{ años; regla del 72: }\SpecialCharTok{\{}\NormalTok{coma(}\DecValTok{72} \OperatorTok{/}\NormalTok{ (}\DecValTok{100} \OperatorTok{*}\NormalTok{ r), }\DecValTok{4}\NormalTok{)}\SpecialCharTok{\}}\SpecialStringTok{ años"}\NormalTok{)}
\end{Highlighting}
\end{Shaded}

\begin{verbatim}
t* = 14,2067 años; regla del 72: 14,4000 años
\end{verbatim}

\begin{atlasnota}

\atlasnotatitulo{Convención de mercado (\mdseries Tasa efectiva anual)}Salvo
indicación contraria, una tasa anual expresada en porcentaje es
efectiva: el capital crece en el factor \(1 + r\) por año. Las
convenciones de capitalización y de conteo de días se estudian en L5
c.~2.

\end{atlasnota}

La llamada \emph{regla del 72} aproxima \(t^{*}\) por \(72 / (100 r)\).
La Tabla~\ref{tbl-regla-72} compara ambos valores; la diferencia
relativa crece con la tasa.

\begin{Shaded}
\begin{Highlighting}[]
\NormalTok{tasas }\OperatorTok{=}\NormalTok{ [}\FloatTok{0.01}\NormalTok{, }\FloatTok{0.02}\NormalTok{, }\FloatTok{0.05}\NormalTok{, }\FloatTok{0.08}\NormalTok{, }\FloatTok{0.10}\NormalTok{]}
\NormalTok{filas }\OperatorTok{=}\NormalTok{ [}\StringTok{"| Tasa | Exacto | Regla del 72 | Diferencia |"}\NormalTok{, }\StringTok{"|{-}{-}{-}:|{-}{-}{-}:|{-}{-}{-}:|{-}{-}{-}:|"}\NormalTok{]}
\ControlFlowTok{for}\NormalTok{ r }\KeywordTok{in}\NormalTok{ tasas:}
\NormalTok{    exacto }\OperatorTok{=}\NormalTok{ math.log(}\DecValTok{2}\NormalTok{) }\OperatorTok{/}\NormalTok{ math.log(}\DecValTok{1} \OperatorTok{+}\NormalTok{ r)}
\NormalTok{    aprox }\OperatorTok{=} \DecValTok{72} \OperatorTok{/}\NormalTok{ (}\DecValTok{100} \OperatorTok{*}\NormalTok{ r)}
\NormalTok{    filas.append(}\SpecialStringTok{f"| }\SpecialCharTok{\{}\NormalTok{coma(}\DecValTok{100} \OperatorTok{*}\NormalTok{ r, }\DecValTok{0}\NormalTok{)}\SpecialCharTok{\}}\SpecialStringTok{ \% | }\SpecialCharTok{\{}\NormalTok{coma(exacto, }\DecValTok{3}\NormalTok{)}\SpecialCharTok{\}}\SpecialStringTok{ | }\SpecialCharTok{\{}\NormalTok{coma(aprox, }\DecValTok{3}\NormalTok{)}\SpecialCharTok{\}}\SpecialStringTok{ | }\SpecialCharTok{\{}\NormalTok{coma(aprox }\OperatorTok{{-}}\NormalTok{ exacto, }\DecValTok{3}\NormalTok{)}\SpecialCharTok{\}}\SpecialStringTok{ |"}\NormalTok{)}
\BuiltInTok{print}\NormalTok{(}\StringTok{"}\CharTok{\textbackslash{}n}\StringTok{"}\NormalTok{.join(filas))}
\end{Highlighting}
\end{Shaded}

{

\begin{longtable}[]{@{}rrrr@{}}

\caption{\label{tbl-regla-72}Tiempo de duplicación exacto y regla del 72
(años).}

\tabularnewline

\toprule\noalign{}
Tasa & Exacto & Regla del 72 & Diferencia \\
\midrule\noalign{}
\endhead
\bottomrule\noalign{}
\endlastfoot
1 \% & 69,661 & 72,000 & 2,339 \\
2 \% & 35,003 & 36,000 & 0,997 \\
5 \% & 14,207 & 14,400 & 0,193 \\
8 \% & 9,006 & 9,000 & −0,006 \\
10 \% & 7,273 & 7,200 & −0,073 \\

\end{longtable}

}

La Figura~\ref{fig-logaritmos} muestra \(\log_b x\) para tres bases.
Todas las curvas pasan por \((1; 0)\), como exige la
Definición~\ref{def-logaritmo}.

\begin{Shaded}
\begin{Highlighting}[]
\ImportTok{import}\NormalTok{ numpy }\ImportTok{as}\NormalTok{ np}
\ImportTok{import}\NormalTok{ matplotlib.pyplot }\ImportTok{as}\NormalTok{ plt}
\ImportTok{from}\NormalTok{ matplotlib.ticker }\ImportTok{import}\NormalTok{ FuncFormatter}
\ImportTok{import}\NormalTok{ matplotlib.font\_manager }\ImportTok{as}\NormalTok{ fm}

\CommentTok{\# Tipografía de las figuras: la misma familia que el texto (Libertinus).}
\NormalTok{FUENTES }\OperatorTok{=} \StringTok{"/usr/share/texlive/texmf{-}dist/fonts/opentype/public/libertinus{-}fonts/"}
\ControlFlowTok{for}\NormalTok{ nombre }\KeywordTok{in}\NormalTok{ (}\StringTok{"LibertinusSerif{-}Regular.otf"}\NormalTok{, }\StringTok{"LibertinusSerif{-}Italic.otf"}\NormalTok{):}
\NormalTok{    fm.fontManager.addfont(FUENTES }\OperatorTok{+}\NormalTok{ nombre)}
\NormalTok{plt.rcParams.update(\{}\StringTok{"font.family"}\NormalTok{: }\StringTok{"Libertinus Serif"}\NormalTok{, }\StringTok{"mathtext.fontset"}\NormalTok{: }\StringTok{"stix"}\NormalTok{\})}

\NormalTok{eje\_coma }\OperatorTok{=}\NormalTok{ FuncFormatter(}\KeywordTok{lambda}\NormalTok{ v, \_: }\SpecialStringTok{f"}\SpecialCharTok{\{}\NormalTok{v}\SpecialCharTok{:g\}}\SpecialStringTok{"}\NormalTok{.replace(}\StringTok{"."}\NormalTok{, }\StringTok{","}\NormalTok{).replace(}\StringTok{"{-}"}\NormalTok{, }\StringTok{"}\CharTok{\textbackslash{}u2212}\StringTok{"}\NormalTok{))}
\NormalTok{x }\OperatorTok{=}\NormalTok{ np.linspace(}\FloatTok{0.1}\NormalTok{, }\DecValTok{8}\NormalTok{, }\DecValTok{400}\NormalTok{)}
\NormalTok{fig, ax }\OperatorTok{=}\NormalTok{ plt.subplots(figsize}\OperatorTok{=}\NormalTok{(}\FloatTok{5.5}\NormalTok{, }\FloatTok{3.2}\NormalTok{))}
\ControlFlowTok{for}\NormalTok{ base, nombre }\KeywordTok{in}\NormalTok{ [(}\DecValTok{2}\NormalTok{, }\StringTok{"$b = 2$"}\NormalTok{), (np.e, }\StringTok{"$b = e$"}\NormalTok{), (}\DecValTok{10}\NormalTok{, }\StringTok{"$b = 10$"}\NormalTok{)]:}
\NormalTok{    ax.plot(x, np.log(x) }\OperatorTok{/}\NormalTok{ np.log(base), label}\OperatorTok{=}\NormalTok{nombre)}
\NormalTok{ax.axhline(}\DecValTok{0}\NormalTok{, color}\OperatorTok{=}\StringTok{"0.6"}\NormalTok{, lw}\OperatorTok{=}\FloatTok{0.8}\NormalTok{)}
\NormalTok{ax.xaxis.set\_major\_formatter(eje\_coma)}
\NormalTok{ax.yaxis.set\_major\_formatter(eje\_coma)}
\NormalTok{ax.set\_xlabel(}\StringTok{"$x$"}\NormalTok{)}
\NormalTok{ax.set\_ylabel(}\VerbatimStringTok{r"}\DecValTok{$}\ErrorTok{\textbackslash{}}\VerbatimStringTok{log\_b x}\DecValTok{$}\VerbatimStringTok{"}\NormalTok{)}
\NormalTok{ax.legend(frameon}\OperatorTok{=}\VariableTok{False}\NormalTok{)}
\NormalTok{fig.tight\_layout()}
\NormalTok{plt.show()}
\end{Highlighting}
\end{Shaded}

\begin{figure}[H]

\centering{

\pandocbounded{\includegraphics[keepaspectratio,alt={Tres curvas crecientes que pasan por el punto (1; 0); la de base 2 es la más empinada.}]{cap-a_files/figure-pdf/fig-logaritmos-output-1.pdf}}

}

\caption{\label{fig-logaritmos}Gráfica de \(\log_b x\) para \(b = 2\),
\(e\) y \(10\).}

\end{figure}%

\section{Ejercicios}\label{sec-log-ejercicios}

\begin{exercise}[]\protect\hypertarget{exr-log-1}{}\label{exr-log-1}

Calcule \(\log_{3} 81\) y \(\log_{1/2} 8\) sin calculadora.

\end{exercise}

\begin{exercise}[]\protect\hypertarget{exr-log-2}{}\label{exr-log-2}

Demuestre que \(\log_b (x / y) = \log_b x - \log_b y\) para
\(x, y > 0\).

\end{exercise}

\begin{exercise}[]\protect\hypertarget{exr-log-3}{}\label{exr-log-3}

¿Cuántos años tarda en triplicarse un capital al \(8\) \% efectivo
anual? Dé el resultado con dos decimales.

\end{exercise}

\bookmarksetup{startatroot}

\chapter{La fórmula de Itô}\label{sec-ito}

Este segundo capítulo de muestra prueba la composición de matemática
avanzada: alineaciones de varias líneas, matrices, definiciones por
casos, una simulación con salida verificada y referencias cruzadas al
capítulo anterior. Sigue la exposición de Øksendal
(\citeproc{ref-oksendal2003}{2003}) y Shreve
(\citeproc{ref-shreve2004b}{2004}); los detalles técnicos están en
Karatzas \& Shreve (\citeproc{ref-karatzas1991}{1991}).

\section{Procesos de Itô}\label{sec-procesos-ito}

Sea \(W = (W_t)_{t \geq 0}\) un movimiento
browniano\index{movimiento browniano@movimiento browniano} estándar
sobre un espacio de probabilidad filtrado
\((\Omega, \mathcal{F}, (\mathcal{F}_t)_{t \geq 0}, \mathbb{P})\).

\begin{definition}[Proceso de
Itô]\protect\hypertarget{def-proceso-ito}{}\label{def-proceso-ito}

Un proceso de Itô\index{proceso de itô@proceso de Itô} es un proceso
adaptado \(X\) de la forma
\begin{equation}\protect\phantomsection\label{eq-proceso-ito}{
X_t = X_0 + \int_0^t \mu_s \, \mathrm{d}s + \int_0^t \sigma_s \, \mathrm{d}W_s , \qquad t \geq 0,
}\end{equation} donde \(\mu\) y \(\sigma\) son procesos adaptados con
\(\int_0^t \lvert \mu_s \rvert \, \mathrm{d}s < \infty\) y
\(\int_0^t \sigma_s^{2} \, \mathrm{d}s < \infty\) casi seguramente. En
forma diferencial se escribe
\(\mathrm{d}X_t = \mu_t \, \mathrm{d}t + \sigma_t \, \mathrm{d}W_t\).

\end{definition}

En dimensión dos, un browniano
\(\mathbf{W}_t = (W^{1}_t, W^{2}_t)^{\top}\) con correlación
\(\rho \in [-1, 1]\) tiene matriz de covarianza
\begin{equation}\protect\phantomsection\label{eq-covarianza}{
\operatorname{Cov}(\mathbf{W}_t) = t \begin{pmatrix} 1 & \rho \\ \rho & 1 \end{pmatrix},
}\end{equation} y la función indicadora de un evento \(A\), que se usará
en L18, se define por casos: \[
\mathbf{1}_A(\omega) = \begin{cases} 1, & \omega \in A, \\ 0, & \omega \notin A. \end{cases}
\]

\begin{theorem}[Fórmula de
Itô]\protect\hypertarget{thm-ito}{}\label{thm-ito}

Sea \(X\) un proceso de Itô como en la Ecuación~\ref{eq-proceso-ito} y
sea \(f \in C^{1,2}([0, \infty) \times \mathbb{R})\). Entonces
\(Y_t = f(t, X_t)\) es un proceso de Itô y, casi seguramente para todo
\(t \geq 0\), \begin{equation}\protect\phantomsection\label{eq-ito}{
\begin{aligned}
f(t, X_t) = f(0, X_0)
  &+ \int_0^t \Bigl( \partial_t f(s, X_s) + \mu_s \, \partial_x f(s, X_s) + \tfrac{1}{2} \sigma_s^{2} \, \partial_{xx} f(s, X_s) \Bigr) \mathrm{d}s \\
  &+ \int_0^t \sigma_s \, \partial_x f(s, X_s) \, \mathrm{d}W_s .
\end{aligned}
}\end{equation}

\end{theorem}

\begin{proof}
Esbozo. Se aproxima \(f\) por un desarrollo de Taylor de segundo orden
sobre una partición de \([0, t]\) de malla que tiende a cero. Los
términos de primer orden convergen a las integrales de la
Ecuación~\ref{eq-ito}; el término
\(\tfrac{1}{2} \partial_{xx} f \, (\Delta X)^{2}\) converge a
\(\tfrac{1}{2} \int_0^t \sigma_s^{2} \partial_{xx} f \, \mathrm{d}s\)
porque la variación cuadrática del browniano en \([0, t]\) es \(t\); los
restantes tienden a cero. La demostración completa está en L18 c.~6.
\end{proof}

\begin{atlasnota}

\atlasnotatitulo{Resultado admitido (\mdseries RA-49, Teorema de Donsker)}El
paseo aleatorio reescalado converge en distribución al movimiento
browniano en el espacio de Skorokhod. La convergencia de las
distribuciones finito-dimensionales se demuestra en L18 §3.6; la
convergencia funcional se admite con referencia a Billingsley
(\citeproc{ref-billingsley1999}{1999}).

\end{atlasnota}

\section{Aplicación: el movimiento browniano geométrico}\label{sec-mbg}

\begin{example}[]\protect\hypertarget{exm-mbg}{}\label{exm-mbg}

Sea
\(\mathrm{d}S_t = \mu S_t \, \mathrm{d}t + \sigma S_t \, \mathrm{d}W_t\)
con \(S_0 > 0\). Aplicando el Teorema~\ref{thm-ito} a
\(f(t, x) = \ln x\), y usando que el logaritmo convierte productos en
sumas (Teorema~\ref{thm-log-producto}),
\begin{equation}\protect\phantomsection\label{eq-mbg}{
\mathrm{d} \ln S_t = \Bigl( \mu - \tfrac{1}{2} \sigma^{2} \Bigr) \mathrm{d}t + \sigma \, \mathrm{d}W_t ,
\qquad
S_t = S_0 \exp\Bigl( \bigl(\mu - \tfrac{1}{2}\sigma^{2}\bigr) t + \sigma W_t \Bigr).
}\end{equation} En particular, \(\mathbb{E}[S_t] = S_0 e^{\mu t}\).

\end{example}

La celda siguiente simula \(100\,000\) trayectorias con el esquema de
Euler-Maruyama (\citeproc{ref-glasserman2003}{Glasserman, 2003}) y
compara la media muestral de \(S_T\) con el valor exacto. La semilla
está fijada para que la salida sea reproducible.

\protect\phantomsection\label{simulacion-mbg}
\begin{Shaded}
\begin{Highlighting}[]
\ImportTok{import}\NormalTok{ numpy }\ImportTok{as}\NormalTok{ np}

\KeywordTok{def}\NormalTok{ coma(x, decimales):}
\NormalTok{    texto }\OperatorTok{=} \SpecialStringTok{f"}\SpecialCharTok{\{}\NormalTok{x}\SpecialCharTok{:,.}\NormalTok{\{decimales\}f}\SpecialCharTok{\}}\SpecialStringTok{"}
    \ControlFlowTok{return}\NormalTok{ texto.replace(}\StringTok{","}\NormalTok{, }\StringTok{" "}\NormalTok{).replace(}\StringTok{"."}\NormalTok{, }\StringTok{","}\NormalTok{)}

\NormalTok{rng }\OperatorTok{=}\NormalTok{ np.random.default\_rng(}\DecValTok{20261016}\NormalTok{)}
\NormalTok{S0, mu, sigma, T, pasos, n }\OperatorTok{=} \FloatTok{100.0}\NormalTok{, }\FloatTok{0.08}\NormalTok{, }\FloatTok{0.20}\NormalTok{, }\FloatTok{1.0}\NormalTok{, }\DecValTok{252}\NormalTok{, }\DecValTok{100\_000}
\NormalTok{dt }\OperatorTok{=}\NormalTok{ T }\OperatorTok{/}\NormalTok{ pasos}
\NormalTok{S }\OperatorTok{=}\NormalTok{ np.full(n, S0)}
\ControlFlowTok{for}\NormalTok{ \_ }\KeywordTok{in} \BuiltInTok{range}\NormalTok{(pasos):}
\NormalTok{    S }\OperatorTok{=}\NormalTok{ S }\OperatorTok{+}\NormalTok{ mu }\OperatorTok{*}\NormalTok{ S }\OperatorTok{*}\NormalTok{ dt }\OperatorTok{+}\NormalTok{ sigma }\OperatorTok{*}\NormalTok{ S }\OperatorTok{*}\NormalTok{ np.sqrt(dt) }\OperatorTok{*}\NormalTok{ rng.standard\_normal(n)}
\NormalTok{media, error }\OperatorTok{=}\NormalTok{ S.mean(), S.std(ddof}\OperatorTok{=}\DecValTok{1}\NormalTok{) }\OperatorTok{/}\NormalTok{ np.sqrt(n)}
\BuiltInTok{print}\NormalTok{(}\SpecialStringTok{f"Media simulada de S\_T: }\SpecialCharTok{\{}\NormalTok{coma(media, }\DecValTok{3}\NormalTok{)}\SpecialCharTok{\}}\SpecialStringTok{ (error estándar }\SpecialCharTok{\{}\NormalTok{coma(error, }\DecValTok{3}\NormalTok{)}\SpecialCharTok{\}}\SpecialStringTok{)"}\NormalTok{)}
\BuiltInTok{print}\NormalTok{(}\SpecialStringTok{f"Valor exacto S0·e\^{}(μT): }\SpecialCharTok{\{}\NormalTok{coma(S0 }\OperatorTok{*}\NormalTok{ np.exp(mu }\OperatorTok{*}\NormalTok{ T), }\DecValTok{3}\NormalTok{)}\SpecialCharTok{\}}\SpecialStringTok{"}\NormalTok{)}
\end{Highlighting}
\end{Shaded}

\begin{verbatim}
Media simulada de S_T: 108,269 (error estándar 0,070)
Valor exacto S0·e^(μT): 108,329
\end{verbatim}

\section{Ejercicios}\label{sec-ito-ejercicios}

\begin{exercise}[]\protect\hypertarget{exr-ito-1}{}\label{exr-ito-1}

Use el Teorema~\ref{thm-ito} para calcular \(\mathrm{d}(W_t^{2})\) y
deduzca que \(W_t^{2} - t\) es una martingala.

\end{exercise}

\bookmarksetup{startatroot}

\chapter*{Bibliografía}\label{bibliografuxeda}
\addcontentsline{toc}{chapter}{Bibliografía}

\markboth{Bibliografía}{Bibliografía}

\protect\phantomsection\label{refs}
\begin{CSLReferences}{1}{0}
\bibitem[\citeproctext]{ref-apostol1967}
Apostol, T. M. (1967). \emph{Calculus, Volume 1: One-Variable Calculus,
with an Introduction to Linear Algebra} (2.ª ed.). John Wiley \& Sons.

\bibitem[\citeproctext]{ref-billingsley1999}
Billingsley, P. (1999). \emph{Convergence of Probability Measures} (2.ª
ed.). Wiley.

\bibitem[\citeproctext]{ref-glasserman2003}
Glasserman, P. (2003). \emph{Monte Carlo Methods in Financial
Engineering} (Vol. 53). Springer.

\bibitem[\citeproctext]{ref-karatzas1991}
Karatzas, I., \& Shreve, S. E. (1991). \emph{Brownian Motion and
Stochastic Calculus} (2.ª ed., Vol. 113). Springer.

\bibitem[\citeproctext]{ref-kellison2009}
Kellison, S. G. (2009). \emph{The Theory of Interest} (3.ª ed.).
McGraw-Hill/Irwin.

\bibitem[\citeproctext]{ref-oksendal2003}
Øksendal, B. (2003). \emph{Stochastic Differential Equations: An
Introduction with Applications} (6.ª ed.). Springer.

\bibitem[\citeproctext]{ref-shreve2004b}
Shreve, S. E. (2004). \emph{Stochastic Calculus for Finance II:
Continuous-Time Models}. Springer.

\bibitem[\citeproctext]{ref-spivak2008}
Spivak, M. (2008). \emph{Calculus} (4.ª ed.). Publish or Perish.

\end{CSLReferences}


\backmatter

\printindex


\end{document}
